\documentclass[11pt]{article}
\usepackage[margin=0.88in]{geometry}
\usepackage[T1]{fontenc}
\usepackage{lmodern}
\usepackage{amsmath,amssymb,amsthm,graphicx,microtype}
\usepackage[hidelinks]{hyperref}

\newtheorem{theorem}{Theorem}
\newtheorem{lemma}[theorem]{Lemma}
\newtheorem{proposition}[theorem]{Proposition}
\newcommand{\D}{\mathbb D}
\newcommand{\T}{\mathbb T}
\newcommand{\C}{\mathbb C}
\newcommand{\Pcal}{\mathcal P}
\newcommand{\X}{\mathcal X}
\newcommand{\A}{\mathcal A_0}
\newcommand{\dd}{\mathop{\rm dist}}
\newcommand{\Om}{\Omega_{\rm max}}

\title{Every Boundary \(F_\sigma\) Set Is a Banach Maximal Domain}
\author{Candidate full solution packet\\
\small arXiv:2407.15997, Problem 4.7}
\date{11 August 2026}

\begin{document}
\maketitle

\noindent\textbf{Status.} Candidate full affirmative answer in the Banach
category, subject to expert review.  For every $F_\sigma$ set
$\Gamma\subset\T$ we construct a Banach space of holomorphic functions on
$\D$ satisfying P1--P3 of Mironov--Sampat and having maximal domain precisely
$\D\cup\Gamma$.

\section{Source problem}

Let $\X\subset\operatorname{Hol}(\D)$ be a Hausdorff topological vector
space such that:
\begin{enumerate}
\item[(P1)] the analytic polynomials are dense in $\X$;
\item[(P2)] evaluation at every point of $\D$ is continuous;
\item[(P3)] multiplication by $z$ is continuous on $\X$.
\end{enumerate}
Its maximal domain is
\[
 \Om(\X)=\{w\in\C:p\mapsto p(w)\text{ is continuous on the
 polynomial subspace of }\X\}.
\]
After realizing $\D\cup\Gamma$ for $\Gamma$ both $F_\sigma$ and
$G_\delta$, Mironov and Sampat ask the following on printed page 13 of
\cite{MS}.

\begin{center}
\includegraphics[width=0.96\textwidth]{figures/open_problem_crop.png}
\end{center}

\begin{quote}
\textbf{Problem 4.7.} Does there exist a Hilbert/Banach space $\mathcal H$
satisfying P1--P3 on $\D$ with $\Om=\D\cup\Gamma$ for any given
$F_\sigma$ set $\Gamma\subset\T$?
\end{quote}

The phrase ``Hilbert/Banach'' permits an affirmative Banach-space
construction.  We do not claim here that the space can always be Hilbert.

\section{Proof intuition}

Write $\Gamma$ as a countable union of closed sets $K_j$.  The natural plan
is to impose a norm for each $K_j$ which makes evaluation continuous on
$K_j$, and then take a summable combination of those norms.  The obstruction
is uniformity: for a boundary peak at $w\notin\Gamma$, the distance from $w$
to $K_j$ may tend to zero arbitrarily quickly as $j$ grows.

The key is a local norm whose peak estimates are bounded independently of
that distance.  Let $\A$ be the analytic coefficient space whose Taylor
coefficients form a $c_0$ sequence.  At $\gamma\in\T$, control
\[
 p(\gamma),\qquad
 \Delta_\gamma p(z):=\frac{p(z)-p(\gamma)}{z-\gamma}
\]
in $\C\oplus\A$.  This is weak enough that the divided differences of the
standard peaks are uniformly bounded.  It is also strong enough to be an
honest function norm: the only possible kernel would contain
$1/(\gamma-z)$, whose Taylor coefficients have constant modulus and hence
do not belong to $c_0$.

Away from the peak, the divided-difference coefficients are oscillatory
tails of a binomial distribution.  Abel summation makes their supremum tend
to zero.  Thus every fixed closed layer contributes a quantity tending to
zero, all layers have a common bound, and dominated convergence handles the
summable outer norm.  An $H^2$ term supplies interior evaluation and a
canonical holomorphic limit; the local injectivity just described rules out
invisible completion vectors.

\section{The local coefficient model}

Define
\[
 \A=\left\{g(z)=\sum_{n=0}^{\infty}a_nz^n:(a_n)_{n\ge0}\in c_0\right\},
 \qquad \|g\|_{\A}=\sup_{n\ge0}|a_n|.
 \tag{1}\label{eq:A0}
\]
Every member of $\A$ is holomorphic on $\D$, convergence in $\A$ implies
locally uniform convergence, polynomials are dense, and multiplication by
$z$ is an isometry.

For $\gamma\in\T$ and a polynomial $p$, put
\[
 \|p\|_\gamma
 :=|p(\gamma)|+\|\Delta_\gamma p\|_{\A}.
 \tag{2}\label{eq:local-norm}
\]

\begin{lemma}[No local ghost vectors]\label{lem:local}
The completion of the polynomials for \eqref{eq:local-norm} is isometric to
$\C\oplus_1\A$ through
\[
 p\longmapsto\bigl(p(\gamma),\Delta_\gamma p\bigr),
 \tag{3}\label{eq:pair-map}
\]
and the associated map into $\operatorname{Hol}(\D)$,
\[
 (c,g)\longmapsto c+(z-\gamma)g(z),
 \tag{4}\label{eq:function-map}
\]
is injective.  Moreover,
\[
 \|zp\|_\gamma\le 2\|p\|_\gamma
 \qquad(p\in\Pcal).
 \tag{5}\label{eq:shift-local}
\]
\end{lemma}

\begin{proof}
Finite Taylor truncations in $\A$ show that every pair $(c,g)$ is a limit
of polynomials $c+(z-\gamma)g_m$, so \eqref{eq:pair-map} has dense range.  If
$c+(z-\gamma)g(z)=0$ on $\D$, then
\[
 g(z)=\frac{c}{\gamma-z}
 =\sum_{n=0}^{\infty}c\gamma^{-n-1}z^n.
\]
The coefficients on the right have constant modulus $|c|$.  Since the
coefficients of $g$ tend to zero, $c=0$, and then $g=0$.  This proves
injectivity.

Writing $p=c+(z-\gamma)g$ gives
\[
 zp=\gamma c+(z-\gamma)(c+zg).
\]
The constant has $\A$-norm $|c|$ and $z$ is an isometry on $\A$; hence
\[
 \|zp\|_\gamma
 \le |c|+|c|+\|g\|_{\A}\le2\|p\|_\gamma.
\]
\end{proof}

For a closed $K\subset\T$, with the empty supremum interpreted as zero,
set
\[
 \|p\|_K:=\sup_{\gamma\in K}\|p\|_\gamma.
 \tag{6}\label{eq:K-norm}
\]

\section{Construction of the Banach function space}

Fix an arbitrary $F_\sigma$ set $\Gamma\subset\T$ and closed sets
$K_j\subset\T$ such that
\[
 \Gamma=\bigcup_{j=1}^{\infty}K_j.
 \tag{7}\label{eq:Gamma-layers}
\]
On the polynomials define
\[
 \|p\|_{\X}
 :=\|p\|_{H^2}+\sum_{j=1}^{\infty}2^{-j}\|p\|_{K_j},
 \tag{8}\label{eq:X-norm}
\]
and let $\X$ be the abstract completion.

\begin{proposition}\label{prop:functionspace}
The natural map from $\X$ into $H^2\subset\operatorname{Hol}(\D)$ is
injective.  Consequently $\X$ is a Banach space of holomorphic functions on
$\D$.  It satisfies P1--P3, and
\[
 \D\cup\Gamma\subset\Om(\X).
 \tag{9}\label{eq:easy-inclusion}
\]
\end{proposition}

\begin{proof}
The $H^2$ term gives a continuous map $J:\X\to H^2$.  We prove that
$J$ has zero kernel.  Let $(p_m)$ be Cauchy for \eqref{eq:X-norm} and suppose
$p_m\to0$ in $H^2$, hence locally uniformly on $\D$.

Fix $j$ and $\gamma\in K_j$.  The pairs
$(p_m(\gamma),\Delta_\gamma p_m)$ are Cauchy in $\C\oplus_1\A$.  If their
limit is $(c_\gamma,g_\gamma)$, the identity
\[
 p_m(z)=p_m(\gamma)+(z-\gamma)\Delta_\gamma p_m(z)
\]
and local uniform convergence give
$c_\gamma+(z-\gamma)g_\gamma(z)=0$ on $\D$.  Lemma~\ref{lem:local} forces
$(c_\gamma,g_\gamma)=0$.

Because $(p_m)$ is Cauchy uniformly over $K_j$, this pointwise conclusion is
uniform: given $\varepsilon>0$, choose $m_0$ such that
$\|p_m-p_n\|_{K_j}<\varepsilon$ for $m,n\ge m_0$, and then let $n\to\infty$
in each local norm.  Thus $\|p_m\|_{K_j}\le\varepsilon$ for $m\ge m_0$.
Finally, Cauchyness in the weighted sum and Fatou's lemma yield
\[
 \sum_j2^{-j}\|p_m\|_{K_j}
 \le \liminf_{n\to\infty}
 \sum_j2^{-j}\|p_m-p_n\|_{K_j}\longrightarrow0.
\]
Together with $\|p_m\|_{H^2}\to0$, this proves $p_m\to0$ in $\X$ and hence
injectivity of $J$.

P1 holds by construction.  Since \eqref{eq:X-norm} dominates the $H^2$ norm,
interior evaluations are continuous, proving P2.  Lemma~\ref{lem:local},
\eqref{eq:K-norm}, and the $H^2$ shift isometry give
$\|zp\|_{\X}\le2\|p\|_{\X}$, so the shift extends boundedly and P3 holds.
Finally, if $\gamma\in K_j$, then
\[
 |p(\gamma)|\le\|p\|_{K_j}\le2^j\|p\|_{\X}.
\]
Thus every point of $\Gamma$, as well as every point of $\D$, belongs to the
maximal domain.
\end{proof}

\section{Uniform boundary peaks}

\begin{lemma}\label{lem:peaks}
Fix $w\in\T$ and define
\[
 q_N(z)=\left(\frac{1+\overline wz}{2}\right)^N.
 \tag{10}\label{eq:peak}
\]
Then:
\begin{enumerate}
\item $q_N(w)=1$ and $\|q_N\|_{H^2}\to0$;
\item $\sup_{N\ge1,\,\gamma\in\T}\|q_N\|_\gamma\le2$;
\item if $K\subset\T$ is closed and $w\notin K$, then
      $\|q_N\|_K\to0$.
\end{enumerate}
\end{lemma}

\begin{proof}
Write
\[
 q_N(z)=\sum_{k=0}^{N}a_{N,k}z^k,
 \qquad
 a_{N,k}=2^{-N}\binom Nk\overline w^{,k}.
\]
The standard binomial identity gives
\[
 \|q_N\|_{H^2}^2
 =4^{-N}\sum_{k=0}^{N}\binom Nk^2
 =4^{-N}\binom{2N}{N}\longrightarrow0.
 \tag{11}\label{eq:H2-peak}
\]

Polynomial division gives
\[
 \Delta_\gamma q_N(z)=\sum_{n=0}^{N-1}b_{N,n}(\gamma)z^n,
 \qquad
 b_{N,n}(\gamma)
 =\sum_{k=n+1}^{N}a_{N,k}\gamma^{k-1-n}.
 \tag{12}\label{eq:division}
\]
Since $|a_{N,k}|=2^{-N}\binom Nk$ and those numbers sum to one,
\[
 |b_{N,n}(\gamma)|\le1.
 \tag{13}\label{eq:uniform-bound}
\]
Also $|q_N(\gamma)|\le1$, proving assertion 2.

For assertion 3 put $u=\overline w\gamma$.  Up to a unimodular factor,
$b_{N,n}(\gamma)$ is the binomial tail
\[
 \sum_{k=n+1}^{N}2^{-N}\binom Nk u^k.
 \tag{14}\label{eq:tail}
\]
Abel summation, splitting the unimodal binomial coefficients at their mode,
shows that for $u\ne1$
\[
 \sup_{0\le n<N}|b_{N,n}(\gamma)|
 \le \frac{C}{|1-u|}
       \max_{0\le k\le N}2^{-N}\binom Nk
 \le \frac{C'}{|w-\gamma|\sqrt{N+1}}.
 \tag{15}\label{eq:Abel}
\]
For completeness, the first inequality uses only
$|\sum_{r=s}^{t}u^r|\le2/|1-u|$ on each monotone half; the second is the
elementary central-binomial estimate.

If $K$ is closed and does not contain $w$, then
$\delta=\dd(w,K)>0$.  Equation \eqref{eq:Abel} makes the $\A$ term in
\eqref{eq:local-norm} tend to zero uniformly for $\gamma\in K$.  Moreover,
\[
 \sup_{\gamma\in K}|q_N(\gamma)|
 \le \left(\sup_{\gamma\in K}
       \left|\frac{1+\overline w\gamma}{2}\right|\right)^N
 =\rho_K^N
\]
with $\rho_K<1$.  Hence $\|q_N\|_K\to0$.
\end{proof}

\section{Main theorem}

\begin{theorem}\label{thm:main}
For every $F_\sigma$ set $\Gamma\subset\T$, the Banach space $\X$ obtained
from \eqref{eq:X-norm} satisfies P1--P3 on $\D$ and
\[
 \Om(\X)=\D\cup\Gamma.
 \tag{16}\label{eq:maximal-domain}
\]
Therefore Problem 4.7 of \cite{MS} has an affirmative answer in the Banach
category.
\end{theorem}

\begin{proof}
Proposition~\ref{prop:functionspace} gives the inclusion from right to left
in \eqref{eq:maximal-domain}.  Let $w\in\T\setminus\Gamma$.  Then $w\notin K_j$ for every
$j$.  Lemma~\ref{lem:peaks} gives
\[
 \|q_N\|_{K_j}\longrightarrow0\quad\text{for every fixed }j,
 \qquad
 \|q_N\|_{K_j}\le2\quad\text{for all }N,j.
\]
Dominated convergence for the counting measure with masses $2^{-j}$, along
with \eqref{eq:H2-peak}, yields
\[
 \|q_N\|_{\X}
 =\|q_N\|_{H^2}+\sum_j2^{-j}\|q_N\|_{K_j}\longrightarrow0.
\]
But $q_N(w)=1$, so polynomial evaluation at $w$ is not continuous on $\X$.

It remains to exclude $|w|>1$.  Put $r_N(z)=(z/w)^N$.  Then
$r_N(w)=1$, $\|r_N\|_{H^2}=|w|^{-N}$, and, for every $\gamma\in\T$,
\[
 |r_N(\gamma)|=|w|^{-N},\qquad
 \Delta_\gamma r_N(z)
 =w^{-N}\sum_{n=0}^{N-1}\gamma^{N-1-n}z^n.
\]
Thus $\|r_N\|_\gamma=2|w|^{-N}$ uniformly in $\gamma$, and
$\|r_N\|_{\X}\le3|w|^{-N}\to0$.  Evaluation at $w$ is again unbounded.
This proves that no point outside $\D\cup\Gamma$ lies in $\Om(\X)$ and
completes the proof.
\end{proof}

\section{Verification, scope, and novelty}

The proof is noncomputational.  A finite script supplied with the packet
checks the polynomial-division formula \eqref{eq:division} and the global coefficient
bound \eqref{eq:uniform-bound}; it is only a sanity check.  The essential reviewer checks
are the Abel estimate \eqref{eq:Abel} and the no-ghost argument in
Proposition~\ref{prop:functionspace}.

The result settles the existence question because a Banach space is one of
the two requested categories.  It does not prove that every such maximal
domain has a Hilbert realization.  It also does not address the separate
higher-dimensional question $\Om=\Omega_{\rm env}$ in the source paper.

A bounded novelty search through 11 August 2026 covered the run indexes,
arXiv:2407.15997v2, DOI 10.1090/proc/17344, exact title and problem wording,
author and citation searches, and combinations of ``maximal domain,''
``$F_\sigma$,'' ``Hilbert/Banach,'' and the rational-boundary example.  The
retrieved records still state only the source's $F_\sigma\cap G_\delta$
Hilbert theorem, and no later answer to Problem 4.7 was found.  This is not a
proof of priority, so novelty confidence is moderate.

\medskip
\noindent\textbf{Human-review recommendation.} Verify \eqref{eq:Abel} uniformly
in the tail index and the passage from pointwise local limits to the
supremum and weighted-sum norms.  Subject to those checks, the construction
is a complete affirmative Banach-space solution.

\begin{thebibliography}{9}
\bibitem{MS}
M. Mironov and J. Sampat,
\emph{Jointly cyclic polynomials and maximal domains},
Proc. Amer. Math. Soc. \textbf{153} (2025), 4781--4795;
arXiv:2407.15997v2, DOI: 10.1090/proc/17344.
\end{thebibliography}

\end{document}
